% Shape-Aware Three-Way Merge for RDF Knowledge Graphs
%
% The SECOND paper in this repository. docs/paper/ is the Quipu store paper
% ("A Governed Bitemporal Knowledge Graph Store"); this is a separate work with
% its own build recipe: `just paper-merge`.
%
% Every number traces to benchmark/mergebench/ (Arm A) or benchmark/replay/
% (Arms B and C). Read those directories' BUILD_REPORT.md before quoting one:
% they are the honesty records, and one of the plan's own hypotheses does not
% survive Arm A.
\documentclass[11pt]{article}

\usepackage[margin=1.1in]{geometry}
\usepackage{amsmath}
\usepackage{amssymb}
\usepackage{booktabs}
\usepackage{enumitem}
\usepackage[hidelinks]{hyperref}
\usepackage{lmodern}
\usepackage[T1]{fontenc}
\usepackage{microtype}
\usepackage{xspace}

\newcommand{\sys}{Quipu\xspace}
\newcommand{\op}{\textsc{shape-aware}\xspace}
\newcommand{\ctx}{\textsc{context-merge}\xspace}
\newcommand{\base}{G_{b}}
\newcommand{\ours}{G_{o}}
\newcommand{\theirs}{G_{t}}

\title{\textbf{Shape-Aware Three-Way Merge for RDF Knowledge Graphs}\\
\large The schema defines the conflict: an implementation and evaluation}
\author{Steve Brown\thanks{\textsc{orcid}
\href{https://orcid.org/0009-0009-1720-1785}{\nolinkurl{0009-0009-1720-1785}}.
With implementation and drafting assistance from Claude (Anthropic).
\textbf{Artifact availability:} the merge operator, the divergence benchmark
(\texttt{benchmark/mergebench/}) and the recorded-divergence replay corpus
(\texttt{benchmark/replay/}) behind every number reported here are in the
development repository, \texttt{github.com/scbrown/quipu}. Each result table
names the command that regenerates it.}}
% Pinned, not \today: main.pdf is committed, so a floating date makes every
% rebuild differ from the committed artifact for a reason unrelated to the
% paper. arXiv stamps its own submission date on the PDF regardless.
\date{30 August 2026}

\begin{document}
\maketitle

\begin{abstract}
Version control works for source code because its merge understands the medium:
lines. Knowledge graphs shared between people and software agents are not lines,
and merging their serialisations makes conflicts out of ordering and misses the
ones that matter. We describe and evaluate a three-way merge for RDF in which
the \emph{schema} decides what a conflict is: the merge is set algebra over
canonical triples, and a slot is contended only where a SHACL shape declares
that it can hold at most one value. Multi-valued predicates union; functional
predicates with divergent values are handed to a person. The same shapes graph
then validates the merge result, so the schema that defines conflicts also
audits their resolution.

We do not claim the first three-way merge for RDF. Git-backed RDF versioning
with merge strategies is established prior work, and the closest neighbour ---
Quit Store's Context Merge --- is included here as a baseline that this operator
extends rather than replaces. Our contribution is an implementation in a
production multi-agent store and an evaluation of what schema-derived conflict
semantics buys against that neighbour and against six other strategies.

On a synthetic divergence benchmark over five seeds, the shape-aware operator
raised $33$ conflicts, all true positives, with no false positives, no triples
lost, none fabricated, and no SHACL violations admitted. The context-overlap
baseline detected the \emph{same} $33$ conflicts and asked $845$ questions to do
it. We also report a limit no schema removes: when two sides mint different
names for one entity, the divergence is invisible to any triple-level operator,
including this one. Replaying a recorded production corpus against the shipped
operator --- $105$ repairs a person actually performed, $93$ of them after
excluding $12$ chained pairs, and $939$ real duplicate-value incidents ---
reproduces both the blindness ($0$ of $93$ alias repairs, matching $0$ of $21$
on the synthetic arm) and the benefit, on data nobody generated for the
purpose.
\end{abstract}

\section{Introduction}

A knowledge graph that more than one writer can edit will diverge, and something
has to put it back together. The obvious move is to reuse the tooling that
already solves this for source code: keep the graph in a file, keep the file in
git, and let a line merge reconcile two versions.

That move fails in a specific and instructive way. A line merge understands a
document as an ordered sequence of lines, and a graph is neither ordered nor a
sequence. Two writers who add unrelated facts about the same entity produce
adjacent edits and a textual conflict. Two writers who re-serialise the same
unchanged graph produce a diff of everything. Meanwhile a writer who sets a
second value on a property the ontology says holds exactly one --- the change
that genuinely cannot be reconciled without a decision --- produces a clean,
conflict-free merge, because the two values live on different lines and nothing
in the file says they compete.

The failure is not that line merge is bad at graphs. It is that conflict is a
\emph{semantic} property of the data model, and a line merge is being asked to
infer it from a syntactic one. Our position is that an RDF merge should take
conflict from the place the data model already defines it: the schema.

\paragraph{This paper.} We describe an operator with two halves. Triples are
canonicalised and merged as sets, so serialisation order cannot manufacture a
conflict. A SHACL shapes graph then decides which slots the set algebra is
allowed to settle on its own: where a shape declares \texttt{sh:maxCount 1} and
the two sides carry different values, the operator refuses and records a
decision; everywhere else, it unions. Because the shapes are already present ---
they are what validates writes into the store --- the merge policy is not a
second configuration to keep in step with the ontology. It \emph{is} the
ontology.

\paragraph{What we are not claiming.} Three-way merge over canonicalised RDF is
not new, and we searched before writing. Quit Store~\cite{arndt2019quit} places
RDF under git with merge strategies, including a Context Merge that identifies
conflicts by node overlap and refers them to a person; TerminusDB offers diff and patch over a versioned graph; RDF
Patch and LD Patch standardise change description; RDFC-1.0 gives canonical
labelling for graphs with blank nodes, which we build on rather than beside. To
our knowledge --- and scoped to that search --- what has not been reported is an
evaluation of \emph{SHACL-derived} conflict semantics against those neighbours
in a deployed system. That evaluation is the contribution, and the framing of
the operator as an extension of Context Merge rather than a replacement for it
is deliberate.

\paragraph{Contributions.}
\begin{enumerate}[leftmargin=*,itemsep=2pt]
  \item A three-way merge operator over canonical triples whose conflict
        predicate is derived from SHACL cardinality, with post-merge validation
        against the same shapes, implemented in a production store
        (\S\ref{sec:operator}, \S\ref{sec:impl}).
  \item A divergence benchmark with an oracle recomputable by a reader from
        published data alone, scoring eight strategies under one accounting
        contract, and a result that contradicts one of our own hypotheses
        (\S\ref{sec:synthetic}).
  \item A replay of recorded multi-agent divergence from a production graph
        against the shipped operator, which reproduces the synthetic arm's
        blind spot on real data and quantifies the benefit on a class that
        production surfaced nothing of (\S\ref{sec:replay}).
  \item A stated limit --- alias divergence is invisible to triple-level merge
        --- reported as a measured false-negative column rather than excluded
        from the benchmark (\S\ref{sec:limits}).
\end{enumerate}

\section{Background}
\label{sec:background}

\paragraph{The unit of change.} An RDF graph is a set of triples. Two
serialisations of one graph may differ in every byte --- Turtle groups by
subject, N-Triples does not, and neither fixes an order within a group. A diff
over serialised text therefore reports differences that do not exist in the
data. Taking the triple as the unit of change removes that entire category:
a three-way merge becomes set algebra over $(\base, \ours, \theirs)$, and no
result depends on how any side chose to write its file.

Blank nodes are the exception, since they have no stable name across
serialisations. RDFC-1.0 (formerly URDNA2015) assigns canonical labels and makes
graph-level equality decidable. Our operator sits on top of that substrate; the
benchmark in \S\ref{sec:synthetic} uses ground graphs only, so it does not
exercise blank-node canonicalisation and is not evidence about it.

\paragraph{Where conflict comes from.} Set algebra alone answers only part of
the question. If both sides add a triple, the union holds both. That is correct
for a person's email addresses and wrong for their date of birth, and nothing in
the triples distinguishes the two cases. The distinction is in the schema: a
SHACL shape with \texttt{sh:maxCount 1} on a property path states that the slot
holds at most one value. A merge that unions into such a slot has not resolved
the divergence; it has produced a graph that violates the very constraint the
store enforces on ordinary writes.

This gives a conflict predicate with three properties we care about. It is
\emph{declared}, not inferred from edit patterns. It is \emph{already present},
because the shapes are what the store validates against. And it is
\emph{auditable}, because the same shapes graph can be run over the merge output
--- so the definition of a conflict and the test for a bad resolution are one
artefact rather than two that can drift apart.

\paragraph{Shares.} The operator reconciles \emph{share bundles}: a directory
holding a canonical N-Triples export, the shapes that govern it, and a manifest
carrying a content hash of the graph and the identifier of the share it was
derived from. The parent pointer is what supplies $\base$, so fork-point
anchoring is by content hash rather than by wall-clock time or by a
centrally-issued sequence number. Two peers who have never spoken can therefore
establish a common ancestor from the bundles alone.

\section{The operator}
\label{sec:operator}

Let $\base$, $\ours$ and $\theirs$ be sets of canonical triples, and let $S$ be a
SHACL shapes graph. A \emph{slot} is a pair $(s, p)$ of subject and predicate;
we write $G[s,p] = \{\, o : (s,p,o) \in G \,\}$ for the values a graph carries in
a slot. The slot is the unit of conflict because it is the unit
\texttt{sh:maxCount} is declared over.

\paragraph{Additions and removals.} Each side's change relative to the base is
$A_x = G_x \setminus \base$ and $R_x = \base \setminus G_x$ for $x \in
\{o, t\}$. These are exact: no heuristic decides whether a triple was
``modified'' rather than removed and re-added, because at triple granularity
there is no such distinction to make.

\paragraph{Functional slots.} Let $\mathit{max}_S(s,p)$ be the maximum
cardinality $S$ declares for $p$ on $s$, or $\infty$ where none is declared. A
slot is \emph{contended} when
\[
  \mathit{max}_S(s,p) = 1
  \;\wedge\;
  \ours[s,p] \neq \theirs[s,p]
  \;\wedge\;
  \ours[s,p] \neq \base[s,p]
  \;\wedge\;
  \theirs[s,p] \neq \base[s,p].
\]
Both sides must have moved. If only one side changed a functional slot, the
other side's value is still the base value, there is no disagreement, and the
change applies --- a merge that stopped to ask about it would be charging a
person for a decision nobody needs to make.

\paragraph{The merge.} If any slot is contended, the operator returns
\texttt{conflicts} and writes nothing; each contended slot is reported with its
base, ours and theirs values, and the declared cardinality that made it a
conflict. Otherwise the result is
\[
  G_m \;=\; \bigl( \ours \cup \theirs \bigr) \setminus \bigl( R_o \cup R_t \bigr),
\]
a retraction on either side being honoured against the other side's retained
triples. The transaction records \emph{both} parent share identifiers, so the
merge is itself provenance: the lineage of a merged graph is inspectable in the
same way a merge commit's is.

\paragraph{Two properties worth stating.} The operator never lands a value for a
slot it flagged --- a conflicted merge writes nothing at all --- so it cannot
score well by both refusing a decision and quietly resolving it. And because
$G_m$ is defined by set operations over canonical triples, it is commutative in
$\ours$ and $\theirs$ and independent of serialisation.

\paragraph{Post-merge validation.} $G_m$ is validated against $S$. This is not
belt-and-braces: it is the falsifiable form of the design claim. If conflict
semantics really are derivable from cardinality, then a merge that consulted the
shapes should not be able to produce a graph those same shapes reject. Any
violation in $G_m$ is corruption the operator admitted without charging anyone
for it, and \S\ref{sec:synthetic} reports that column for every strategy
precisely so the claim can fail.

\section{Implementation}
\label{sec:impl}

The operator ships in \sys, a governed bitemporal RDF store, as
\texttt{share\_merge} (approximately 480 lines of Rust). It exposes two entry
points. \texttt{status} loads an incoming bundle, resolves the common ancestor
through the manifest's parent pointer, and reports what each side added and
removed together with the contended slots --- without writing anything, so a
peer can be inspected before it is merged. \texttt{merge} performs the algebra
of \S\ref{sec:operator} and commits, recording both parent share identifiers on
the transaction.

Three implementation choices are worth naming because they affect what the
evaluation can claim.

\paragraph{The shapes come from the bundle, not the store.} Cardinality is read
from the shapes carried inside the share being merged, not from the local
store's copy. A peer's constraints therefore travel with its data, and merging
does not silently reinterpret their graph under our schema.

\paragraph{Conflict is reported at slot granularity, with its evidence.} A
\texttt{DecisionRecord} carries the subject, the predicate, the declared
\texttt{maxCount}, and the base, ours and theirs value sets. The person
resolving it is not told that a conflict exists and left to find it; they are
told which constraint made it one and what the three candidate states were.

\paragraph{The evaluated operator is the shipped one.} The benchmark of
\S\ref{sec:synthetic} contains its own reference implementation of every arm,
including ours, so that all eight are scored under one accounting contract. The
replay of \S\ref{sec:replay} deliberately does \emph{not}: it drives
\texttt{share\_merge::\{status,merge\}} through real bundles built by the
production share path. This matters because a benchmark's copy of an operator
can agree with a benchmark's copy of an oracle for reasons that have nothing to
do with the deployed system.

\section{Arm A: a synthetic divergence benchmark}
\label{sec:synthetic}

\paragraph{Design.} One seeded base graph over a synthetic vocabulary is copied
twice and edited independently --- separate random streams, so neither side's
edits can be a function of the other's --- under a six-operation model: assert,
retract, re-describe, re-type, alias-mint, node-add. An overlap parameter
controls how often an edit targets a contended quarter of the entity set, so
collision probability is set directly rather than left to chance. Eight
strategies then merge $(\base, \ours, \theirs)$ and are scored against one
oracle. There is no language model anywhere in the loop; the writers are
deterministic, so a run is its own ground truth.

\paragraph{The oracle is a function of published data.} \texttt{ground\_truth}
takes the three graphs and the shapes and nothing else --- no edit log, no
generator intent --- so a reader holding only the published inputs recomputes it
exactly, and a generator bug cannot hide inside it. An instrument control
asserts precisely that.

\paragraph{The accounting contract.} Every arm answers two questions whose
answers are in tension. \emph{Human decisions} is the number of slots it refused
to settle; a refused slot is held at its base value in the output, so no arm can
flag a conflict \emph{and} land a value for it. \emph{SHACL violations} counts
violations in what it landed automatically, against the same shapes that defined
the conflicts --- corruption it admitted without charging anyone. Either can be
driven to zero by sacrificing the other, so neither is reportable alone.
\emph{Triples lost} and \emph{triples spurious} are the third leg, and they
catch the arm that admits no corruption and asks no questions because it
discarded one side's work.

\paragraph{Results.} Table~\ref{tab:armA} aggregates five seeds
($1, 7, 42, 99, 2026$) at $200$ entities, $400$ edits per side and overlap
$0.5$. Regenerate with \texttt{just bench merge --seed <N>}; instrument controls
with \texttt{just bench merge --selftest} ($7$ properties, non-zero exit on any
failure).

\begin{table}[t]
\centering\small
\begin{tabular}{lrrrrrr}
\toprule
arm & decisions & precision & recall & SHACL & lost & spurious \\
\midrule
\texttt{git-turtle-reserialized} & 681--794 & 0.006--0.011 & 0.36--0.80 & 0--1 & 266--299 & 62--71 \\
\texttt{git-turtle-stable}       & 41--59   & 0.049--0.119 & 0.20--0.70 & 0--3 & 39--57   & 12--24 \\
\texttt{git-canonical}           & 106--144 & 0.028--0.057 & 0.30--0.70 & 0--3 & 95--128  & 18--31 \\
\texttt{union}                   & 0        & ---          & 0.00       & 53--72 & 2--7   & 112--136 \\
\texttt{lww-theirs}              & 0        & ---          & 0.00       & 0--1 & 7--16    & 4--10 \\
\texttt{triple-3way}             & 0        & ---          & 0.00       & 3--6 & 2--7     & 7--17 \\
\ctx                             & 159--182 & 0.025--0.046 & 0.50--0.80 & 0    & 128--150 & 63--72 \\
\op                              & 4--8     & \textbf{1.000} & 0.50--0.80 & \textbf{0} & \textbf{0} & \textbf{0} \\
\bottomrule
\end{tabular}
\caption{Eight strategies, five seeds, ranges across seeds. Recall is bounded
above by $1$ minus the alias-mint fraction for \emph{every} arm
(\S\ref{sec:limits}). The line-merge arms shell out to real \texttt{git
merge-file}; where git is absent they report \emph{unavailable}, never zero.}
\label{tab:armA}
\end{table}

Three readings matter.

\paragraph{The comparison with the nearest neighbour is not about detection.}
Summed over the five seeds the oracle contains $54$ conflicts, of which $21$ are
alias-mint and $33$ are visible at triple level. \op raised exactly those $33$,
with zero false positives at every individual seed. \ctx raised the
\emph{same} $33$ --- identical recall, seed by seed --- and $812$ others, for
$845$ questions in total. The schema does not find conflicts the node-overlap
heuristic misses. It distinguishes the ones that need a person from the $812$
that do not, which is a $25.6\times$ difference in the quantity the operator is
actually spending: attention.

\paragraph{The zero columns are not all worth the same.} \texttt{union},
\texttt{lww-theirs} and \texttt{triple-3way} each report $0$ human decisions, and
none of them resolved anything. \texttt{union} admitted $53$--$72$ SHACL
violations. \texttt{lww-theirs} admitted almost none --- and silently dropped
$7$--$16$ triples and fabricated $4$--$10$, which is what last-writer-wins looks
like when you count what it discarded rather than what it complained about. Only
\op holds all three at zero while still asking.

\paragraph{A hypothesis of ours did not survive.} We predicted that
canonicalisation alone would materially beat non-canonical line merge. Against
the re-serialised arm it does: unparseable output lines go from $116$--$169$ to
$0$. Against a \emph{stably} serialised Turtle file it does not.
\texttt{git-canonical} raises \emph{more} conflicts ($106$--$144$) than
\texttt{git-turtle-stable} ($41$--$59$), because sorting scatters each subject's
triples into one sorted run where two sides' additions land adjacent and collide,
whereas subject-grouped Turtle absorbed them into a block. The honest statement
is narrower than the hypothesis: canonicalisation buys syntactic integrity and
independence from serialisation order, and it does not by itself reduce conflict
count. It is a trade, and we report it as one.

\paragraph{The circularity, stated rather than managed away.} \op and the oracle
both read the same shapes graph, and for the triple-visible classes their rules
coincide by construction. \op's precision of $1.000$ on this arm is therefore a
property of the design, not evidence about it. We do not repair this by making
the oracle cleverer, which would only move the coincidence somewhere less
visible. We handle it by reporting per-class results so a reader can see which
rows are definitional, by confining this arm's contribution to the \emph{cost of
the alternatives}, and by taking correctness evidence from \S\ref{sec:replay},
where the operator under test is the shipped one and disagreement is a finding.

\section{Arms B and C: recorded divergence, shipped operator}
\label{sec:replay}

The synthetic arm can measure what the alternatives cost. It cannot be evidence
that our operator is correct, because its oracle and its operator share a
definition. This arm removes that circularity twice over: the corpus is
\emph{recorded} production divergence rather than generated, and the operator
under test is the shipped \texttt{share\_merge} driven through real bundles.
Here a disagreement between operator and expectation is a finding rather than a
harness bug.

\paragraph{The corpus.} Extracted from a live multi-agent knowledge graph in
which software agents and people had been writing concurrently for months, and
committed as a fixture. Entity names are replaced by random labels and the
mapping is discarded at build time; structure, cardinality and class membership
are preserved exactly, and injectivity is verified over the $1{,}076$ names.
Two divergence classes were recorded.

\emph{Alias-mint}: two names silently minted for one entity, later repaired by a
person asserting \texttt{owl:sameAs}. There are $171$ raw \texttt{sameAs} edges,
which reduce to $105$ distinct undirected repairs ($66$ edges were one repair
recorded twice). Each is one repair a person actually performed --- the
manual-effort baseline. \textbf{The class splits, and the split matters}: $52$
of the $105$ are \emph{id-form}, two spellings of the same commit identifier,
which a normalisation rule could retire without asking anyone; the other $53$
are two English phrasings of one concept, where nothing but a person can decide.
Reporting $105$ as one undifferentiated ``human decisions'' figure would
overstate the irreducible cost roughly twofold.

\emph{Comment-double}: a \texttt{sh:maxCount 1} description predicate that the
production append path doubled instead of replacing --- $939$ affected subjects
carrying $1{,}784$ excess values.

\paragraph{Results.} Table~\ref{tab:armB} reports the replay. Regenerate with
\texttt{cargo run --example replay --features shacl}.

\begin{table}[t]
\centering\small
\begin{tabular}{lrrl}
\toprule
scenario & decisions & historical & outcome \\
\midrule
\texttt{alias-id-form}   & 0   & 46 & merged --- \textbf{false negative} \\
\texttt{alias-semantic}  & 0   & 47 & merged --- \textbf{false negative} \\
\texttt{comment-double}  & 939 & 0  & conflicts --- raised as decisions \\
\texttt{sameas-repair}   & 0   & 20 & merged --- repair survives (wanted) \\
\bottomrule
\end{tabular}
\caption{Recorded divergence replayed against the shipped operator. $12$ pairs
($6$ id-form, $6$ semantic) were excluded as \emph{chained} --- sharing an
endpoint with another pair --- because including them measures an incidental
collision on the shared endpoint rather than alias detection.}
\label{tab:armB}
\end{table}

\paragraph{The operator is blind to the entire alias class: $0$ of $93$.} This
reproduces the synthetic arm's $0$ of $4$ on data nobody generated for the
purpose, and it is the honest false-negative column. Two names for one entity is
invisible to \emph{any} triple-level operator, shape-aware included; no schema
fixes it, because it is a limit of the medium the operator works in.

\paragraph{The operator raises $939$ decisions on a class production surfaced
none of.} The same class produced $1{,}784$ silent excess values through the
append path. This is a counterfactual and must be read as one
(\S\ref{sec:limits}).

\paragraph{A repair survives a concurrent edit.} \texttt{owl:sameAs} carries no
cardinality constraint, so a repair asserted on one side unions rather than
racing an edit to the aliased node on the other. A merge that could drop repairs
would make the repair path itself unreliable, so this outcome is the wanted one
rather than an incidental one.

\paragraph{The negative control.} Rerunning with the description predicate's
\texttt{maxCount} unbound takes \texttt{comment-double} from $939$ decisions to
$0$, everything merging cleanly. The decisions are therefore \emph{shape-derived}
--- not manufactured by the harness --- and this control arm doubles as the
closest thing to a simulation of the production append path: it reproduces the
doubling. That is the paper's thesis stated as an experiment. Unbind the
constraint and the conflict ceases to exist; the schema is what made it one.

\paragraph{Arm C: share, diverge, reconnect.} A case study walks the full
lifecycle on a real bundle and verifies all eight hops: share a base bundle;
a peer publishes a divergent share naming it as parent; we edit locally; status
before reconnect reports \texttt{diverged} with $+40$ triples each side and $0$
conflicts; the merge asserts $40$ and retracts $0$; the transaction records two
parents; both sides' work is present afterwards ($40/40$ each); and status after
reconnect reports nothing outstanding. One detail is worth carrying: the graph
hash is stable across runs while the share identifier is not, because the
manifest embeds a timestamp. Fork-point anchoring depends on the graph hash, so
lineage is unaffected --- but a share identifier must not be used as a content
identity.

\section{Limitations}
\label{sec:limits}

\paragraph{Alias divergence is invisible, and that is a property of the medium.}
When two writers mint different names for one entity, the graphs disagree about
the world while agreeing about every triple. No triple-level operator can see
it, so a recall of $1.000$ is not achievable by anything in Table~\ref{tab:armA}.
We left the class in the benchmark deliberately: it is the most common real
divergence defect in the recorded corpus, and excluding it would have produced a
shape-aware arm scoring $1.000$ on both precision and recall as an artefact of
what we chose to measure. Resolving it needs entity resolution --- a different
kind of system, working from labels and neighbourhoods rather than from triples,
and one whose errors are not obviously cheaper than the ones it prevents.

\paragraph{The $939$ is a counterfactual.} Those doublings arose from sequential
re-posts through a single write path, not from a share, a divergence and a merge.
The supported claim is narrow: had the same edits arrived as a divergence and
been reconciled by this operator, each would have been surfaced as a decision
instead of silently doubling. It is \emph{not} a claim that the operator was
running and prevented $1{,}784$ doublings.

\paragraph{Comment bodies in the corpus are synthetic.} Real values are prose
about a deployment. Only the multiplicity is real --- which is all the merge
reasons about, but it does mean the corpus cannot support any claim about value
content.

\paragraph{A conflict blocks the whole merge, not the conflicting slot.} The
operator returns \texttt{conflicts} with nothing asserted. On a slice like the
$939$-subject one, a single reconnect writes nothing until every contended slot
is resolved. This is the conservative choice and it is a real usability cost;
partial application with a held-back conflict set is future work.

\paragraph{Ground graphs only.} The benchmark's graphs contain no blank nodes,
so canonicalisation there reduces to sorting. RDFC-1.0 is the substrate we build
on and is not exercised by this arm; nothing here is evidence about blank-node
canonicalisation.

\paragraph{Timing is indicative.} One machine, one process, single runs. The
line-merge arms include process spawn for \texttt{git merge-file}, of order
$5$\,ms, which dominates their timing at every size measured; those columns are
not comparable with the in-process arms and we do not report them as if they
were.

\paragraph{Slot granularity is an assumption.} A conflict is reported over a
$(subject, predicate)$ slot because that is the unit \texttt{sh:maxCount} is
declared over. An operator working at another granularity would need a different
scorer, and the comparison in Table~\ref{tab:armA} would not transfer unchanged.

\paragraph{One deployment.} The recorded corpus comes from a single production
store. Its class distribution --- alias-mint and functional-value doubling
dominating --- may not generalise, and we make no claim about divergence rates
in other systems from it.

\section{Related work}
\label{sec:related}

\paragraph{Git-backed RDF versioning.} Quit Store~\cite{arndt2019quit} is the
closest neighbour and the reason this paper makes no priority claim. It places
RDF under git, versions quads with commit-level provenance, and supplies merge
strategies --- including a Context Merge that reconciles by node overlap. The
triple-set framing of three-way merge is present there; we are not introducing
it.

Our delta is where the conflict predicate comes from. Context Merge derives
contention from the \emph{shape of the change} --- whether the two sides touched
overlapping context. We derive it from the \emph{shape of the data}, as declared
in SHACL. The consequence is measurable rather than rhetorical, and it is the
central result of \S\ref{sec:synthetic}: across five seeds the two approaches
detect an identical set of $33$ triple-visible conflicts, and the context
heuristic asks $845$ questions where cardinality asks $33$. Overlap is a good
proxy for ``these edits are related''. It is a poor proxy for ``these edits
cannot both be true'', and only the second needs a person. We therefore position
this operator as an extension of that line of work: the algebra is theirs, the
conflict predicate is the contribution.

\paragraph{Diff and patch over versioned graphs.} TerminusDB~\cite{terminusdb} provides diff and
patch over a versioned graph store. RDF Patch and LD Patch standardise the
description of change sets. These solve the representation and transport of a
delta; they are complementary to, and do not decide, what makes two deltas
irreconcilable.

\paragraph{Canonicalisation.} W3C RDF Dataset Canonicalization~\cite{rdfc2024} (RDFC-1.0,
formerly URDNA2015) gives deterministic labelling of blank nodes and hence
decidable graph equality. We build directly on it: it is what makes ``the same
triple'' well defined across two independently serialised bundles, and it is a
substrate rather than a competitor. R43ples~\cite{graube2014r43ples} and archiving systems such as
OSTRICH address efficient storage and querying of revision histories, a
different problem from reconciling two live divergent copies.

\paragraph{Convergent replicated types.} CRDT constructions for RDF, such as
SU-Set~\cite{ibanez2012suset}, guarantee convergence without coordination. They and this operator make
opposite trades deliberately. A CRDT resolves every case automatically, which
requires that some rule always picks a winner; for a functional predicate whose
two values are both asserted in good faith, any such rule silently discards
somebody's knowledge. We would rather stop and ask, and accept that this means
a merge can fail. Where automatic convergence is required, a CRDT is the right
tool and this is not; where the graph is a shared record of what is believed to
be true, a silently dropped assertion is the more expensive failure. Operational
transformation has the same shape of trade in a different lineage.

\paragraph{Schema-aware merging elsewhere.} Ontology alignment and matching
address a related question --- reconciling different conceptualisations --- but
operate on schema-to-schema correspondence rather than on three-way
reconciliation of instance data under a shared schema. Structured merge tools
for source code (semi-structured and AST merges) share our motivation exactly:
raise the merge above lines to the level at which the medium defines conflict.
Our claim is the RDF instance of that principle, with SHACL supplying the
structure.

\section{Conclusion}

Merging knowledge graphs by merging their text asks a line-oriented tool to
infer a semantic property, and it fails in both directions: it manufactures
conflicts from serialisation order and misses the divergences that genuinely
cannot be reconciled automatically. Taking the canonical triple as the unit of
change removes the first failure entirely. Taking the conflict predicate from a
SHACL shapes graph addresses the second, and does so without adding a policy
that has to be maintained alongside the ontology --- because it \emph{is} the
ontology, and the same shapes then audit the result.

The evaluation supports a narrow claim, which is the one we make. Against the
nearest neighbour, schema-derived conflict semantics do not detect more: they
detect the same $33$ triple-visible conflicts across five seeds while asking
$33$ questions rather than $845$, with no triples lost, none fabricated and no
constraint violations admitted. Replaying $939$ real duplicate-value incidents and
$93$ of $105$ recorded repairs ($12$ chained pairs excluded) against the shipped
operator reproduces that benefit on data nobody generated for the purpose --- and reproduces the blind spot too:
alias divergence is invisible to triple-level merge, $0$ of $93$, and no schema
will fix it.

Two directions follow. Partial application would let a merge land its
uncontended slots and hold back only the contended ones, removing the cost
noted in \S\ref{sec:limits}. And the alias class, which dominates the recorded
corpus and which this operator cannot see, is where entity resolution would have
to meet merge --- with the open question being whether its errors are cheaper
than the divergence it repairs.


\bibliographystyle{plain}
\bibliography{references}

\end{document}
